% ------------------------------------------------------------------------------------
% formulation.tex — the score-bundle model, stated once, self-consistently.
%
% Canonical notation (must match CLAUDE.md and the code):
%   S = {s_i}, s_i = (p_i, b_i, d_i)   score support (pitch, beat onset, beat duration)
%   y_i = [tau_i, log r_i, v_i]        Phase-1 per-note variables
%   Q_G                                graph prior precision (lambda I + eta L_G, or Matern)
%   lambda, eta                        additive ridge term, Laplacian weight
%   sigma_g, kappa, alpha              Matern scale, inverse-length, exponent
%   Sigma_e / Sigma_y / m              observation-noise cov / posterior cov / posterior mean
%   sigma                              POSTERIOR standard deviation only
%   mu_LM, h_i                         LM prior mean, LM per-note embedding
%   z, a, Phi(z), x, eps, A_i(t)       Phase-3 positions, amplitudes, synth, audio, noise,
%                                      amplitude envelope
% ------------------------------------------------------------------------------------
\documentclass[11pt]{article}
\usepackage[margin=1.1in]{geometry}
\usepackage{amsmath,amssymb,bm}
\usepackage{booktabs}

\newcommand{\R}{\mathbb{R}}
\newcommand{\N}{\mathcal{N}}
\newcommand{\score}{\mathcal{S}}
\newcommand{\graph}{\mathcal{G}}
\newcommand{\Qg}{Q_{\graph}}
\newcommand{\Lg}{L_{\graph}}
\newcommand{\muLM}{\mu_{\mathrm{LM}}}
\newcommand{\diag}{\operatorname{diag}}

\title{Score-Bundle Models: Formulation\\
\large Bayesian score-informed performance transcription with a graph prior}
\author{}
\date{}

\begin{document}
\maketitle

\section{Setting}

A \emph{score} is the known discrete support
\begin{equation}
  \score = \{ s_i \}_{i=1}^{N}, \qquad s_i = (p_i,\, b_i,\, d_i),
  \label{eq:support}
\end{equation}
where $p_i$ is the pitch (MIDI number), $b_i \ge 0$ the onset in score time (beats), and
$d_i > 0$ the notated duration (beats). A \emph{performance} attaches to each score event a
vector of continuous expressive variables; in Phase 1 (piano),
\begin{equation}
  y_i \;=\; \bigl[\, \tau_i,\; \log r_i,\; v_i \,\bigr] \in \R^{3},
  \label{eq:variables}
\end{equation}
with
\begin{align}
  \tau_i    &= t_i - T(b_i)                       && \text{onset residual (seconds): performed onset $t_i$ minus tempo-implied onset,} \nonumber\\
  \log r_i  &= \log \frac{\delta_i}{T(b_i + d_i) - T(b_i)} && \text{articulation: log ratio of performed duration $\delta_i$ to warp-implied duration,} \nonumber\\
  v_i       &= \tilde v_i / 127 - \overline{\tilde v/127}  && \text{dynamics: centred, unit-range MIDI velocity $\tilde v_i$.} \nonumber
\end{align}
Here $T(\cdot)$ is the monotone time warp (score beats $\to$ performance seconds) estimated
from the alignment's beat grids; conditioning on $T$ is what makes the per-note field $y$
well defined and jointly Gaussian. Later phases add channels on the same support:
intonation $c_i$ (cents), vibrato parameters, and the fundamental curve $f_0$ (Phase~2);
harmonic amplitudes $a$ (Phase~3). The task is \emph{score-informed}: $\score$ is given,
and we infer the posterior over $y$ (inverse direction) from observations generated by the
same model that, run forward, synthesizes a performance.

\section{The score graph}

Let $\graph = (V, E)$ be a weighted graph on the $N$ score events with Gaussian-kernel
edge weights
\begin{equation}
  W_{ij} \;=\; \exp\!\Bigl( -\tfrac{(b_i - b_j)^2}{2\,\ell_b^2} - \tfrac{(p_i - p_j)^2}{2\,\ell_p^2} \Bigr)
  \;+\; w_v\, \mathbf{1}[\text{voice}_i = \text{voice}_j],
  \qquad W_{ii} = 0,
  \label{eq:adjacency}
\end{equation}
with length scales $\ell_b$ (beats) and $\ell_p$ (semitones), an optional same-voice bonus
$w_v \ge 0$, and optional $k$-nearest-neighbour sparsification. The (combinatorial) graph
Laplacian is
\begin{equation}
  \Lg \;=\; D - W, \qquad D = \diag\Bigl(\textstyle\sum_j W_{ij}\Bigr).
  \label{eq:laplacian}
\end{equation}
The \emph{temporal-only} baseline replaces $W$ by the chain adjacency of score-time order
(each note coupled only to its temporal neighbours), an AR(1)-style smoother; the
\emph{independent} baseline drops all edges ($\Lg = 0$).

\section{Graph prior over performance variables}

Each channel of $y$ (a scalar field over the $N$ nodes) receives a Gaussian Markov random
field prior whose precision is built from $\Lg$ in one of two interchangeable forms:
\begin{align}
  \text{additive:} \qquad & \Qg \;=\; \lambda I + \eta\, \Lg,
      && \lambda > 0 \text{ (ridge / pull to the mean)}, \;\; \eta \ge 0 \text{ (coupling weight)};
  \label{eq:additive}\\[2pt]
  \text{Mat\'ern / SPDE:} \qquad & \Qg \;=\; \sigma_g^{-2}\,\bigl(\kappa^2 I + \Lg\bigr)^{\alpha},
      && \sigma_g > 0 \text{ (scale)}, \;\; \kappa > 0 \text{ (inverse length)}, \;\; \alpha \in \mathbb{Z}_{\ge 1} \text{ (exponent)}.
  \label{eq:matern}
\end{align}
For integer $\alpha$ both forms are sparse GMRF precisions. The joint prior over the
stacked field ($k$ channels) is block-diagonal across channels, each with its own
hyperparameters. Throughout, $\alpha$ is \emph{only} the Mat\'ern exponent, $\eta$ is
\emph{only} the additive Laplacian weight, and $\sigma_g$ is the Mat\'ern prior scale ---
plain $\sigma$ is reserved for the posterior standard deviation
(\S\ref{sec:posterior}).

\paragraph{Prior mean.}
The field is centred on a mean $\mu$ from one of three sources: $\mu = 0$; a hand-built
ridge regression on local score features; or the \emph{learned} mean $\muLM$. For the
latter, a pretrained autoregressive music language model supplies a per-note embedding
$h_i \in \R^{d}$ (its hidden state at note $i$'s last token), and a ridge head
$W_{\mathrm{head}}$, fit on a disjoint split, maps embeddings to means:
\begin{equation}
  \mu_{\mathrm{LM},i} \;=\; W_{\mathrm{head}}^{\top} \bigl[ h_i;\, 1 \bigr],
  \qquad\quad
  y \;=\; \muLM + \xi, \qquad \xi \sim \N\bigl(0,\, \Qg^{-1}\bigr).
  \label{eq:lm-mean}
\end{equation}
The LM supplies the expressive signal; the graph prior models the \emph{residual} with
correlated, calibrated uncertainty. This turns the core comparison into: does a
score-graph residual prior on top of a learned mean improve recovery \emph{and}
calibration over the mean alone, and over hand-built baselines?

\section{Observation model and posterior}
\label{sec:posterior}

Observations are the latent field plus independent Gaussian noise,
\begin{equation}
  \tilde y \;=\; y + e, \qquad e \sim \N(0,\, \Sigma_e),
  \qquad \Sigma_e = \diag(\sigma_{e,1}^2, \dots, \sigma_{e,N}^2)
  \;\;\text{(homoscedastic } \Sigma_e = \sigma_e^2 I \text{ unless stated)}.
  \label{eq:obs}
\end{equation}
Masked (held-out) notes contribute no likelihood term: writing $P = \Sigma_e^{-1}$ on
observed nodes and $0$ on masked nodes (a diagonal observation-precision matrix), the
posterior over the whole field is closed-form Gaussian,
\begin{equation}
  p(y \mid \tilde y) = \N\bigl(m,\, \Sigma_y\bigr), \qquad
  \Sigma_y \;=\; \bigl(\Qg + P\bigr)^{-1}, \qquad
  m \;=\; \Sigma_y \bigl(\Qg\, \mu + P\, \tilde y\bigr).
  \label{eq:posterior}
\end{equation}
Masked notes are predicted from their graph neighbours through the conditional Gaussian
--- the held-out \emph{imputation} task. The per-note posterior standard deviation is
$\sigma_i = \sqrt{(\Sigma_y)_{ii}}$.

\paragraph{Predictive variance for held-out observations (the variance floor).}
A held-out \emph{observation} at node $i$ is $\tilde y_i = y_i + e_i$, so its predictive
variance must include the observation noise:
\begin{equation}
  \operatorname{Var}\bigl[\tilde y_i \mid \tilde y_{\mathrm{obs}}\bigr]
  \;=\; (\Sigma_y)_{ii} + \sigma_e^2 .
  \label{eq:floor}
\end{equation}
The latent variance $(\Sigma_y)_{ii}$ collapses toward zero at well-pinned nodes, so
scoring held-out data against $(\Sigma_y)_{ii}$ alone is catastrophically overconfident
(NLL and coverage blow up); Eq.~\eqref{eq:floor} is required for all held-out evaluation.

\section{Hyperparameters (empirical Bayes)}
\label{sec:eb}

With the additive form \eqref{eq:additive}, the per-piece, per-channel hyperparameters
$\theta = (\lambda, \eta, \sigma_e^2)$ are chosen by maximizing the marginal likelihood of
the observed nodes (indices $o$),
\begin{equation}
  \log p(\tilde y_o \mid \theta)
  \;=\; \log \N\bigl(\tilde y_o;\; \mu_o,\; K_{oo} + (\Sigma_e)_{oo}\bigr),
  \qquad K = \Qg^{-1},
  \label{eq:marglik}
\end{equation}
optimized in log-parameter space (positivity) with a derivative-free simplex method. A
\emph{calibration-split} alternative instead splits the observed nodes, and picks $\theta$
minimizing the held-out predictive NLL of the calibration subset under the floored
predictive variance \eqref{eq:floor}.

\paragraph{Noise floor.}
On a minority of mask realizations the unconstrained maximizer of \eqref{eq:marglik}
degenerates ($\sigma_e^2 \to 0$), producing an overconfident posterior; the same principle
as \eqref{eq:floor} applies to the \emph{fit}: constrain $\sigma_e^2 \ge \varepsilon_{\mathrm{nv}}$
(a small floor, e.g.\ a fixed fraction of the observed residual variance) inside the
optimization.

\section{Evaluation: recovery and calibration}

Predictions on held-out notes are scored on both axes. With held-out values $y_i$,
predictive means $m_i$ and floored predictive standard deviations
$\sigma_i^{\mathrm{pred}} = \sqrt{(\Sigma_y)_{ii} + \sigma_e^2}$
(over $n$ held-out notes):
\begin{align}
  \text{RMSE} &= \Bigl( \tfrac1n \textstyle\sum_i (y_i - m_i)^2 \Bigr)^{1/2},
  &
  \text{NLL} &= \tfrac1n \textstyle\sum_i \Bigl[ \tfrac12 \log\bigl(2\pi\, (\sigma_i^{\mathrm{pred}})^2\bigr)
        + \tfrac{(y_i - m_i)^2}{2 (\sigma_i^{\mathrm{pred}})^2} \Bigr],
  \\[4pt]
  \text{cov}@\gamma &= \tfrac1n \textstyle\sum_i \mathbf{1}\bigl[\, |y_i - m_i| \le z_{\gamma}\, \sigma_i^{\mathrm{pred}} \bigr],
  &
  \text{cal-err} &= \sup_u \bigl|\, \hat F_{\mathrm{PIT}}(u) - u \,\bigr|,
\end{align}
where $z_\gamma$ is the two-sided standard-normal quantile for level $\gamma$ (0.9
throughout) and $\hat F_{\mathrm{PIT}}$ is the empirical CDF of the probability integral
transform $u_i = \Phi_{\N}\bigl((y_i - m_i)/\sigma_i^{\mathrm{pred}}\bigr)$, which is
Uniform$(0,1)$ iff the predictive distribution is calibrated
($\Phi_{\N}$: standard normal CDF --- distinct from the Phase-3 synthesis operator
$\Phi(z)$). A win for the structured prior requires improving \emph{both} recovery (RMSE)
and calibration (NLL, coverage, cal-err) against the independent, temporal-only, ridge,
and LM-mean-only baselines, on held-out, contamination-filtered data.

\section{Phase 2: intonation and vibrato (continuously pitched instruments)}

For voice/strings/winds, new channels on the same support: intonation
\begin{equation}
  c_i \;=\; 1200 \, \log_2 \frac{f_0^{(i)}}{f_{\mathrm{ref}}\, 2^{\,p_i / 12}}
  \qquad \text{(cents, relative to the equal-tempered target)},
\end{equation}
plus vibrato descriptors (rate Hz, extent cents, onset delay) extracted from the per-note
$f_0$ curve. These channels reuse the graph prior and posterior
\eqref{eq:additive}--\eqref{eq:posterior} unchanged; only the targets and graph length
scales differ.

\section{Phase 3: waveform likelihood}

The optional audio likelihood replaces $\tilde y$ with the waveform $x \in \R^{M}$:
\begin{equation}
  x \;=\; \Phi(z)\, a + \epsilon, \qquad \epsilon \sim \N(0,\, K_n),
  \label{eq:waveform}
\end{equation}
where $z$ collects the \emph{nonlinear position variables} (tempo, timing, intonation,
vibrato --- the Phase-1/2 fields), $\Phi(z) \in \R^{M \times 2K}$ is the deterministic
harmonic synthesis basis (quadrature columns $\cos\phi_k(t), \sin\phi_k(t)$ with
cumulative phase $\phi_k(t) = 2\pi k \int_0^t f_0(s)\, ds$ determined by $z$), and
$a \in \R^{2K}$ are the harmonic amplitudes. Per-note amplitudes are shaped in time by an
amplitude envelope $A_i(t)$ (attack--decay support of note $i$), which multiplies the
columns of $\Phi(z)$ belonging to note $i$. With a Gaussian amplitude prior
$a \sim \N(\mu_a, \Sigma_a)$, the amplitudes marginalize exactly:
\begin{align}
  p(a \mid x, z) &= \N(m_a,\, S_a), &
  S_a &= \bigl( \Sigma_a^{-1} + \Phi^\top K_n^{-1} \Phi \bigr)^{-1}, &
  m_a &= S_a \bigl( \Sigma_a^{-1} \mu_a + \Phi^\top K_n^{-1} x \bigr),
  \label{eq:amp-post}\\[2pt]
  p(x \mid z) &= \N\bigl( x;\; \Phi\, \mu_a,\; \Phi\, \Sigma_a\, \Phi^\top + K_n \bigr),
  \label{eq:collapsed}
\end{align}
writing $\Phi = \Phi(z)$. Inference over the nonlinear $z$ (Laplace approximation or
variational inference, initialized from the nominal score) is the open Phase-3 step; the
graph prior \eqref{eq:additive} plays the role of $p(z)$'s structured component.

\section{Notation summary}

\begin{center}
\begin{tabular}{ll}
\toprule
Symbol & Meaning \\
\midrule
$\score = \{s_i\}$, $s_i = (p_i, b_i, d_i)$ & score support: pitch, beat onset, beat duration \\
$y_i = [\tau_i, \log r_i, v_i]$ & Phase-1 variables: onset residual, articulation, dynamics \\
$T(\cdot)$ & tempo warp, score beats $\to$ performance seconds \\
$\graph$, $W$, $\Lg$ & score graph, adjacency \eqref{eq:adjacency}, Laplacian \eqref{eq:laplacian} \\
$\ell_b$, $\ell_p$, $w_v$ & graph length scales (beats, semitones), voice bonus \\
$\Qg$ & graph prior precision: $\lambda I + \eta \Lg$ or $\sigma_g^{-2}(\kappa^2 I + \Lg)^\alpha$ \\
$\lambda$, $\eta$ & additive ridge term, Laplacian weight \\
$\sigma_g$, $\kappa$, $\alpha$ & Mat\'ern scale, inverse length, exponent \\
$\mu$, $\muLM$, $h_i$ & prior mean; LM-predicted mean; LM per-note embedding \\
$\tilde y$, $e$, $\Sigma_e$ & observations, observation noise, its covariance \\
$m$, $\Sigma_y$, $\sigma_i$ & posterior mean, posterior covariance, posterior std (only) \\
$\sigma_i^{\mathrm{pred}}$ & floored predictive std, $\sqrt{(\Sigma_y)_{ii} + \sigma_e^2}$ \\
$c_i$, $f_0$ & intonation (cents), fundamental frequency (Phase 2) \\
$x$, $z$, $a$, $\Phi(z)$, $\epsilon$, $K_n$ & audio, nonlinear positions, amplitudes, synth basis, noise, noise cov (Phase 3) \\
$A_i(t)$ & amplitude envelope of note $i$ (Phase 3) \\
\bottomrule
\end{tabular}
\end{center}

\paragraph{Deliberate disambiguations.}
$\score$ (support) is never a covariance; plain $\sigma$ is never a prior scale ($\sigma_g$
is); $\alpha$ is never the additive weight ($\eta$ is); $A_i(t)$ (capital) is the amplitude
envelope, never lowercase $a_i(t)$; $\Phi(z)$ is the synthesis operator, $\Phi_{\N}$ the
normal CDF.

\end{document}
